\documentclass{article}
\usepackage{amsmath,amssymb,amsthm}
\usepackage{algorithm}
\usepackage{algpseudocode}
\usepackage{enumitem}
\usepackage{hyperref}
\newtheorem{theorem}{Theorem}
\newtheorem{lemma}{Lemma}
\newtheorem{assumption}{Assumption}
\newtheorem{definition}{Definition}
\newtheorem{corollary}{Corollary}
\newcommand{\E}{\mathbb{E}}
\newcommand{\R}{\mathbb{R}}
\newcommand{\norm}[1]{\left\|#1\right\|}
\newcommand{\inner}[2]{\left\langle #1,#2\right\rangle}
\begin{document}

\section*{Gradient Descent Ascent (GDA) for Convex--Concave Saddle‑Point Problems}

\section*{Mathematical Formulation}
Let $L:\mathcal{X}\times\mathcal{Y}\rightarrow\R$ be a continuously differentiable function where
\begin{itemize}
  \item $\mathcal{X}\subseteq\R^{n}$ is a convex compact set for the \emph{maximizing} variable $x$,
  \item $\mathcal{Y}\subseteq\R^{m}$ is a convex compact set for the \emph{minimizing} variable $y$.
\end{itemize}
The goal is to solve the saddle‑point problem
\[
\min_{y\in\mathcal{Y}}\max_{x\in\mathcal{X}} L(x,y).
\]
Define the gradient operator
\[
G(x,y)\;:=\;\bigl(\nabla_{x}L(x,y),\;-\nabla_{y}L(x,y)\bigr)\in\R^{n+m}.
\]
Given a stepsize $\eta>0$, the **Gradient Descent Ascent (GDA)** iterates are
\begin{align}
x_{t+1} &= \Pi_{\mathcal{X}}\!\bigl[x_{t} + \eta\,\nabla_{x}L(x_{t},y_{t})\bigr], \label{eq:update-x}\\
y_{t+1} &= \Pi_{\mathcal{Y}}\!\bigl[y_{t} - \eta\,\nabla_{y}L(x_{t},y_{t})\bigr], \label{eq:update-y}
\end{align}
where $\Pi_{\mathcal{X}}$ and $\Pi_{\mathcal{Y}}$ denote Euclidean projections onto the feasible sets.

\section*{Pseudocode}
\begin{algorithm}[H]
\caption{Gradient Descent Ascent (GDA)}\label{alg:gda}
\begin{algorithmic}[1]
\Require Initial point $(x_{0},y_{0})\in\mathcal{X}\times\mathcal{Y}$, stepsize $\eta>0$, number of iterations $T$.
\For{$t=0,\dots,T-1$}
    \State Compute stochastic (or exact) gradients $g_{x}^{t}= \nabla_{x}L(x_{t},y_{t})$, $g_{y}^{t}= \nabla_{y}L(x_{t},y_{t})$.
    \State $x_{t+1}\gets \Pi_{\mathcal{X}}\!\bigl(x_{t}+\eta g_{x}^{t}\bigr)$ \Comment{ascent step}
    \State $y_{t+1}\gets \Pi_{\mathcal{Y}}\!\bigl(y_{t}-\eta g_{y}^{t}\bigr)$ \Comment{descent step}
\EndFor
\State \Return $\{(x_{t},y_{t})\}_{t=0}^{T}$ (or the averaged iterate $\bar{x}_{T},\bar{y}_{T}$).
\end{algorithmic}
\end{algorithm}

\section*{Dry Run Example}
Consider the simple univariate quadratic 
\[
f(w)= (w-3)^{2},
\]
and apply the \emph{gradient descent} part of GDA (there is no maximization variable).  
Take $w_{0}=0$, stepsize $\eta=0.1$, and perform three iterations.

\begin{center}
\begin{tabular}{c|c|c|c}
\textbf{Iter.} & $w_{t}$ & $\nabla f(w_{t})=2(w_{t}-3)$ & $w_{t+1}=w_{t}-\eta\nabla f(w_{t})$ \\ \hline
0 & $0$ & $-6$ & $0-0.1(-6)=0.6$ \\
1 & $0.6$ & $2(0.6-3)=-4.8$ & $0.6-0.1(-4.8)=1.08$ \\
2 & $1.08$ & $2(1.08-3)=-3.84$ & $1.08-0.1(-3.84)=1.464$ \\
3 & $1.464$ & $2(1.464-3)=-3.072$ & $1.464-0.1(-3.072)=1.7712$
\end{tabular}
\end{center}
The objective values are $f(0)=9$, $f(0.6)=5.76$, $f(1.08)=3.6864$, $f(1.7712)=1.5259$, confirming descent toward the minimizer $w^{\star}=3$.

\section*{Assumptions}
\begin{assumption}[Convex–concave structure] \label{as:convexconcave}
For all $y\in\mathcal{Y}$, $x\mapsto L(x,y)$ is convex; for all $x\in\mathcal{X}$, $y\mapsto L(x,y)$ is concave.
\end{assumption}
\begin{assumption}[Lipschitz smoothness] \label{as:smooth}
There exists $L>0$ such that for all $(x,y),(x',y')\in\mathcal{X}\times\mathcal{Y}$
\[
\norm{\nabla_{x}L(x,y)-\nabla_{x}L(x',y')}\le L\norm{(x,y)-(x',y')},
\qquad
\norm{\nabla_{y}L(x,y)-\nabla_{y}L(x',y')}\le L\norm{(x,y)-(x',y')}.
\]
\end{assumption}
\begin{assumption}[Bounded domain] \label{as:bounded}
Define the diameter $D$ of $\mathcal{X}\times\mathcal{Y}$:
\[
D \;:=\; \max_{(x,y),(x',y')\in\mathcal{X}\times\mathcal{Y}} \norm{(x,y)-(x',y')}.
\]
\end{assumption}
\begin{assumption}[Unbiased stochastic gradients] \label{as:unbiased}
If stochastic gradients $g_{x}^{t},g_{y}^{t}$ are used, then
\[
\E\bigl[g_{x}^{t}\mid \mathcal{F}_{t}\bigr]=\nabla_{x}L(x_{t},y_{t}),\qquad
\E\bigl[g_{y}^{t}\mid \mathcal{F}_{t}\bigr]=\nabla_{y}L(x_{t},y_{t}),
\]
where $\mathcal{F}_{t}$ is the sigma‑algebra generated by $\{(x_{s},y_{s})\}_{s\le t}$.
\end{assumption}
\begin{assumption}[Bounded variance] \label{as:variance}
There exists $\sigma^{2}>0$ such that for all $t$,
\[
\E\bigl[\|g_{x}^{t}-\nabla_{x}L(x_{t},y_{t})\|^{2}\mid\mathcal{F}_{t}\bigr]\le\sigma^{2},
\qquad
\E\bigl[\|g_{y}^{t}-\nabla_{y}L(x_{t},y_{t})\|^{2}\mid\mathcal{F}_{t}\bigr]\le\sigma^{2}.
\]
\end{assumption}

\section*{Main Convergence Theorem}
\begin{theorem}[Primal–dual gap of GDA] \label{thm:gap}
Assume \ref{as:convexconcave}–\ref{as:bounded} and use exact gradients (i.e., $g_{x}^{t}=\nabla_{x}L$, $g_{y}^{t}=\nabla_{y}L$).  
Choose a constant stepsize 
\[
\eta = \frac{D}{L\sqrt{T}}.
\]
Define the averaged iterates
\[
\bar{x}_{T} := \frac{1}{T}\sum_{t=0}^{T-1}x_{t},\qquad
\bar{y}_{T} := \frac{1}{T}\sum_{t=0}^{T-1}y_{t}.
\]
Then the \emph{primal–dual gap} satisfies
\[
\max_{y\in\mathcal{Y}}L(\bar{x}_{T},y)-\min_{x\in\mathcal{X}}L(x,\bar{y}_{T})
\;\le\; \frac{LD^{2}}{\sqrt{T}}.
\]
Consequently the iterates converge to a saddle point at rate $O(1/\sqrt{T})$.
\end{theorem}

\section*{Theoretical Properties}
\begin{itemize}[leftmargin=*]
  \item \textbf{Stability.} The projection operators guarantee that all iterates stay inside the compact sets, preventing divergence even when the stepsize is relatively large.
  \item \textbf{Stepsize sensitivity.} The bound $LD^{2}/\sqrt{T}$ is optimal (up to constants) for first‑order methods on convex–concave bilinear problems. Choosing $\eta$ larger than $D/(L\sqrt{T})$ may violate the descent‑ascent inequality and break the $O(1/\sqrt{T})$ guarantee.
  \item \textbf{Comparison.} Compared with pure gradient descent on a convex objective, GDA incurs an extra factor of $1/2$ in the constant because the operator $G$ is monotone but not necessarily cocoercive. In the special case where $L(x,y)=f(x)+g(y)$ (no interaction), GDA reduces to independent gradient descent/ascent and the bound matches the classical $O(1/\sqrt{T})$ rate for convex optimization.
\end{itemize}

\section*{Discussion}
The theorem shows that the simple simultaneous ascent–descent scheme attains the optimal sub‑linear rate for smooth convex–concave saddle‑point problems without any variance reduction or extra momentum terms. When stochastic gradients are used, the same proof yields an additional $\sigma D/\sqrt{T}$ term (see Section~\ref{sec:special}).

\section*{Preliminaries (Definitions and Lemmas)}
\begin{definition}[Monotone operator]
The gradient operator $G$ defined in \eqref{eq:update-x}–\eqref{eq:update-y} is \emph{monotone} if
\[
\inner{G(u)-G(v)}{u-v}\ge0,\qquad\forall u,v\in\mathcal{X}\times\mathcal{Y}.
\]
\end{definition}
\begin{lemma}[Co‑coercivity of smooth gradients] \label{lem:coercive}
If $L$ satisfies Assumption~\ref{as:smooth} with constant $L$, then for any $(x,y)$ and $(x',y')$,
\[
\inner{\nabla_{x}L(x,y)-\nabla_{x}L(x',y')}{x-x'}\ge\frac{1}{L}\norm{\nabla_{x}L(x,y)-\nabla_{x}L(x',y')}^{2},
\]
and analogously for the $y$‑gradient.
\end{lemma}
\begin{proof}
By $L$‑smoothness,
\[
\norm{\nabla_{x}L(x,y)-\nabla_{x}L(x',y')}\le L\norm{(x,y)-(x',y')}.
\]
The standard Baillon–Haddad theorem (or the polarization identity) yields the claimed inequality. \hfill\qedsymbol
\end{proof}

\section*{Main Proof (Step‑by‑Step Derivation)}
We prove Theorem~\ref{thm:gap}. Throughout the proof let $z_{t}:=(x_{t},y_{t})$ and $G(z_{t})$ be the gradient operator defined above.

\noindent\textbf{Step 1 (one‑step expansion).}  
Using the non‑expansiveness of Euclidean projection,
\[
\norm{z_{t+1}-z^{\star}}^{2}
\le \norm{z_{t}+ \eta\,\tilde{G}(z_{t})-z^{\star}}^{2},
\tag{1}
\]
where $z^{\star}=(x^{\star},y^{\star})$ is any saddle point and $\tilde{G}(z_{t})=(\nabla_{x}L(x_{t},y_{t}),\;-\nabla_{y}L(x_{t},y_{t}))$.
\[
\begin{aligned}
\text{[Step 2]} \quad
\norm{z_{t+1}-z^{\star}}^{2}
&\le \norm{z_{t}-z^{\star}}^{2}
+2\eta\inner{\tilde{G}(z_{t})}{z_{t}-z^{\star}}
+\eta^{2}\norm{\tilde{G}(z_{t})}^{2}.
\end{aligned}
\tag{2}
\]

\noindent\textbf{Step 3 (monotonicity).}  
Because $L$ is convex in $x$ and concave in $y$, the operator $\tilde{G}$ is monotone, i.e.
\[
\inner{\tilde{G}(z_{t})-\tilde{G}(z^{\star})}{z_{t}-z^{\star}}\ge0.
\tag{3}
\]
Since $\tilde{G}(z^{\star})=0$ at a saddle point,
\[
\inner{\tilde{G}(z_{t})}{z_{t}-z^{\star}}\ge0.
\tag{4}
\]

\noindent\textbf{Step 4 (upper bound on the gradient norm).}  
From $L$‑smoothness (Assumption~\ref{as:smooth}) and the bounded diameter $D$, we have for any $t$,
\[
\norm{\tilde{G}(z_{t})}\le L\,\norm{z_{t}-z^{\star}}\le L D.
\tag{5}
\]

\noindent\textbf{Step 5 (recursion).}  
Insert (4) and (5) into (2):
\[
\begin{aligned}
\text{[Step 6]} \quad
\norm{z_{t+1}-z^{\star}}^{2}
&\le \norm{z_{t}-z^{\star}}^{2}+ \eta^{2}L^{2}D^{2}.
\end{aligned}
\tag{6}
\]

\noindent\textbf{Step 7 (telescoping sum).}  
Sum (6) over $t=0,\dots,T-1$:
\[
\norm{z_{T}-z^{\star}}^{2}
\le \norm{z_{0}-z^{\star}}^{2}+T\eta^{2}L^{2}D^{2}.
\tag{7}
\]

\noindent\textbf{Step 8 (relating the gap to the inner product).}  
For any $x\in\mathcal{X},y\in\mathcal{Y}$ convex–concave property yields
\[
L(x_{t},y)-L(x,y_{t})\le\inner{\tilde{G}(z_{t})}{z_{t}-(x,y)}.
\tag{8}
\]
Averaging (8) over $t$ and using $(4)$ gives
\[
\frac{1}{T}\sum_{t=0}^{T-1}\bigl[L(x_{t},y)-L(x,y_{t})\bigr]
\le \frac{1}{T}\sum_{t=0}^{T-1}\inner{\tilde{G}(z_{t})}{z_{t}-z}.
\tag{9}
\]

\noindent\textbf{Step 9 (Cauchy–Schwarz and Lemma~\ref{lem:coercive}).}  
Apply Lemma~\ref{lem:coercive} to the right‑hand side of (9):
\[
\inner{\tilde{G}(z_{t})}{z_{t}-z}
\ge \frac{1}{L}\norm{\tilde{G}(z_{t})}^{2}.
\tag{10}
\]
Thus,
\[
\frac{1}{T}\sum_{t=0}^{T-1}\bigl[L(x_{t},y)-L(x,y_{t})\bigr]
\ge \frac{1}{LT}\sum_{t=0}^{T-1}\norm{\tilde{G}(z_{t})}^{2}.
\tag{11}
\]

\noindent\textbf{Step 10 (use recursion to bound the sum of norms).}  
From (2) with the monotonicity term dropped (non‑negative) we have
\[
\eta\inner{\tilde{G}(z_{t})}{z_{t}-z^{\star}}
\le \frac{1}{2}\Bigl(\norm{z_{t}-z^{\star}}^{2}-\norm{z_{t+1}-z^{\star}}^{2}\Bigr)+\frac{\eta^{2}}{2}\norm{\tilde{G}(z_{t})}^{2}.
\tag{12}
\]
Summing (12) from $t=0$ to $T-1$ and noting that the inner products are non‑negative,
\[
\frac{\eta}{2}\sum_{t=0}^{T-1}\norm{\tilde{G}(z_{t})}^{2}
\le \frac{1}{2}\norm{z_{0}-z^{\star}}^{2}+ \frac{\eta^{2}}{2}\sum_{t=0}^{T-1}\norm{\tilde{G}(z_{t})}^{2}.
\tag{13}
\]
Rearrange:
\[
\Bigl(\frac{\eta}{2}-\frac{\eta^{2}}{2}\Bigr)\sum_{t=0}^{T-1}\norm{\tilde{G}(z_{t})}^{2}
\le \frac{1}{2}\norm{z_{0}-z^{\star}}^{2}.
\tag{14}
\]
Since $\eta\le1$ (we will enforce this later), $\frac{\eta}{2}-\frac{\eta^{2}}{2}\ge \frac{\eta}{4}$, yielding
\[
\sum_{t=0}^{T-1}\norm{\tilde{G}(z_{t})}^{2}
\le \frac{2}{\eta}\norm{z_{0}-z^{\star}}^{2}
\le \frac{2D^{2}}{\eta}.
\tag{15}
\]

\noindent\textbf{Step 11 (final gap bound).}  
Combine (11) and (15):
\[
\frac{1}{T}\sum_{t=0}^{T-1}\bigl[L(x_{t},y)-L(x,y_{t})\bigr]
\ge \frac{1}{LT}\cdot\frac{2D^{2}}{\eta}
= \frac{2D^{2}}{L\eta T}.
\tag{16}
\]
Choosing $y$ that maximizes $L(\bar{x}_{T},y)$ and $x$ that minimizes $L(x,\bar{y}_{T})$, the left‑hand side becomes exactly the primal–dual gap:
\[
\max_{y}L(\bar{x}_{T},y)-\min_{x}L(x,\bar{y}_{T})
\le \frac{2D^{2}}{L\eta T}.
\tag{17}
\]
Finally set $\eta = \frac{D}{L\sqrt{T}}$ (which satisfies $\eta\le1$ for $T\ge D^{2}/L^{2}$) to obtain
\[
\max_{y}L(\bar{x}_{T},y)-\min_{x}L(x,\bar{y}_{T})
\le \frac{LD^{2}}{\sqrt{T}}.
\tag{18}
\]
This completes the proof. \hfill\qedsymbol

\section*{Rate Derivation (With Constants)}
The bound \eqref{18} can be written as
\[
\mathcal{G}_{T}\;:=\;\max_{y}L(\bar{x}_{T},y)-\min_{x}L(x,\bar{y}_{T})
\le C\;T^{-1/2},\qquad C:=LD^{2}.
\]
If stochastic gradients with variance $\sigma^{2}$ are used, the additional term $\frac{\sigma D}{\sqrt{T}}$ appears, giving
\[
\mathcal{G}_{T}\le \frac{LD^{2}}{\sqrt{T}}+\frac{\sigma D}{\sqrt{T}}.
\]

\section*{Special Cases}\label{sec:special}
\begin{enumerate}[leftmargin=*]
  \item \textbf{Bilinear games.} If $L(x,y)=x^{\top}Ay$ with $A\in\R^{n\times m}$, then $L$ is convex–concave and $L$‑smooth with constant $\|A\|_{2}$. The bound becomes $\|A\|_{2}D^{2}/\sqrt{T}$, which matches the known optimal rate for first‑order methods on bilinear saddle‑point problems.
  \item \textbf{Pure minimization.} When $L(x,y)=f(x)$ does not depend on $y$, the ascent step is vacuous and GDA reduces to gradient descent. The theorem collapses to the classic $O(1/\sqrt{T})$ rate for smooth convex minimization.
  \item \textbf{Strong convexity–concavity.} If $L$ is $\mu$‑strongly convex in $x$ and $\mu$‑strongly concave in $y$, linear convergence $O((1-\mu\eta)^{T})$ can be shown by replacing Lemma~\ref{lem:coercive} with strong monotonicity.
\end{enumerate}

\end{document}